\appendix
\renewcommand{\theHsection}{\Alph{section}}
\renewcommand{\theHtable}{\Alph{section}.\arabic{table}}
\renewcommand{\theHfigure}{\Alph{section}.\arabic{figure}}

\section{Exponents at a conical point}
\setcounter{table}{0}
\label{app:cone}

The tails at the vertices of the cube and at the tip of the almond need the exponents of the field at a conical point. Near the vertex of a perfectly conducting cone the total field is a sum of terms of two types \cite{vanbladel1996,desmedt1986,costabel2000},
\begin{equation}
  \bE = \nabla\bigl(r^{\nu}\,Y(\hat\bx)\bigr), \qquad \bH = \nabla\bigl(r^{\tau}\,Z(\hat\bx)\bigr),
\end{equation}
where $Y$ and $Z$ are eigenfunctions of the spherical Laplace--Beltrami operator (the angular part of the Laplacian on the unit sphere) on the exterior spherical region outside the cone. They satisfy the Dirichlet condition $Y = 0$ and the Neumann condition $\partial_n Z = 0$ along the cone boundary, with eigenvalues $\nu(\nu+1)$ and $\tau(\tau+1)$. The electric type gives $\bE\sim r^{\nu-1}$, whereas the magnetic type gives $\bH\sim r^{\tau-1}$ and $\bE\sim r^{\tau}$. For polyhedral corners these are the principal singular parts of the time-harmonic field found by Costabel and Dauge \cite[Section 6.6]{costabel2000}.

We compute these eigenvalues by the subspace-angle method of particular solutions \cite{betcke2005}, which avoids spurious near-zero solutions by normalising the trial space at interior sample points on the sphere. For a trial value of $\nu$, the Legendre functions $P_\nu(-\hat\bx\cdot\mathbf{s}_j)$ with source poles $\mathbf{s}_j$ on the unit sphere inside the cone satisfy the spherical eigenvalue equation $(\Delta_{S^2}+\nu(\nu+1))u = 0$ away from their logarithmic singularity at $\hat\bx = \mathbf{s}_j$. Since they are the fundamental solutions on the sphere, this is again an MFS, and $\nu$ is an exponent whenever an interior-normalised combination vanishes on the boundary curve. At the almond tip, the poles $\mathbf{s}_j$ lie on the focal segment of the elliptic cross-section in $z = 0$ and cluster towards its ends, with the normal derivative $\partial_{s_z}P_\nu$ used for modes odd in $z$. For the cube corner, they lie inside the spherical triangle that the cube cuts out and cluster exponentially towards its three corners, using $P_\nu$ together with its two tangential derivatives at each pole. As a check on a circular cone, whose exponent is known as a root of a Legendre function, the method reproduces it to eight digits.

\emph{Cube vertex.} For the exterior of a corner of the cube we obtain $\nu = 0.45417$ and $1.23087$, and $\tau = 0.840$ and $1.20635$. Here we give the first magnetic exponent to three digits because the Neumann singularity at the corners of the spherical triangle converges more slowly. De Smedt and Van Bladel \cite{desmedt1986} give $\nu = 0.45410$ and $\tau = 0.84108$ for the rectangular corner.

\emph{Almond tip.} The tip is an elliptic cone with half-angles $34.09^\circ$ and $12.71^\circ$, that is, with axis ratio $1/3$. \autoref{tab:coneexp} lists the exponents by the parities of the potential under the two mirrors of the almond. Only one exponent is below $0.87$, namely $\nu = 0.29959$, which belongs to the class that is even under both mirrors. Because $\bH\propto\nabla\times\bE$ is an axial vector (it acquires an extra sign change under mirror reflection), the magnetic scalar potential of a parity class has the opposite mirror parities to its electric potential. Thus, the class that is even under both mirrors has the electric exponents of the first row of the table and the magnetic exponents of the last row. The values are unchanged in five digits when the number of poles is varied from $12$ to $36$. Bilinear interpolation in the tables of \cite{desmedt1986} gives $\nu = 0.298$ and $\tau = 0.885$ and $0.967$ for this cone, against our $0.29959$, $0.87905$ and $0.96803$. The differences of up to $0.7\%$ are of the size of the accuracy stated there.

\begin{table}[htbp]
  \centering
  \caption{Exponents of the elliptic cone at the tip of the almond, by the parities of the potential under the mirrors $y\to-y$ and $z\to-z$. The two smallest exponents of each type are listed.}
  \label{tab:coneexp}
  \begin{tabular}{lll}
    \toprule
    Parity $(y,z)$ & $\nu$ (electric type) & $\tau$ (magnetic type) \\
    \midrule
    even, even & $0.29959$, $1.46427$ & $1.05604$, $1.98580$ \\
    odd, even  & $1.10372$, $2.24730$ & $0.96803$, $1.96545$ \\
    even, odd  & $1.04120$, $2.10576$ & $0.87905$, $1.75454$ \\
    odd, odd   & $2.01295$, $3.04849$ & $1.97538$, $2.90094$ \\
    \bottomrule
  \end{tabular}
\end{table}

\emph{From exponents to weights.} The tail sources lie on rows that run into the conical point with a spacing proportional to the distance $r$. At short range, electric dipoles of strength $r^{e}$ on such a row produce an electric field $\bE\sim r^{e-3}$ (since $\bG\sim r^{-3}$), whereas magnetic dipoles of strength $r^{e}$ produce $\bH\sim r^{e-3}$ and, by \eqref{eq:magnetic}, $\bE\sim r^{e-2}$. Matching $\bE\sim r^{\nu-1}$ therefore gives $e = \nu+2$ for electric dipoles, and matching $\bH\sim r^{\tau-1}$ (and $\bE\sim r^{\tau}$) gives $e = \tau+2$ for magnetic dipoles. Note that letting each type be carried by the dipoles of its own kind is the simplest assignment, and we have not derived it. For the cube, a vertex tail with two exponents for each kind of dipole changes the error at $J = 6$ by less than $1\%$.

\section{Evaluation of the line sources}
\setcounter{table}{0}
\label{app:lines}

A line source of the cube at depth $d$ carries a polynomial density in the stretched variable $u$, with $n = n_0 + \lceil 2\ln(1/d)\rceil$ Legendre modes for each of its three components, where $n_0 = \max(4, 2J-4)$. At two and three wavelengths each line has four and twelve more modes (\autoref{sec:cubeother}). Each matrix entry is the integral of a dipole field against one mode. The integrand varies on the scale of the distance between the target and the line, which ranges from a fraction of $d$ to the size of the cube. We organise the quadrature so that most of the work is shared between the targets.

The interval $-1\le u\le1$ is divided into a binary tree of panels, and each panel carries a fixed set of nodes. On the four coarsest levels, a panel near the vertices is wide relative to the stretched variable $u(t)$, so the modes $P_m(u(t))$ vary rapidly in $t$. Because these panels are used only for distant targets, the Green's function kernel is smooth in $t$ across them. There we sample the smooth kernel at $12$ Gauss--Legendre nodes in the physical variable $t$ and precompute the exact integrals (moments) of the corresponding Lagrange interpolating polynomials against the modes $P_m(u(t))$. On the finer levels we use $8$ Gauss nodes in $u$. All nodes and weights are computed once for each depth. A target descends the tree and accepts a panel as soon as the length of the panel is smaller than its distance from the target. A distant target then needs $12$ kernel values for the whole line, whereas a target at distance $d$ needs a few hundred.

The electric and the magnetic line at the same position share their kernel values. Likewise, the nine lines of a tail share the quadrature nodes and mode weights of the line at $d_{\min}$ and are summed before multiplying by the modes. Because each line is symmetric across the plane through its midpoint, only four of the eight mirror images of a target need to be evaluated, and their signed sums give the matrices of all parity classes in one pass.

With $20$ and $10$ nodes in place of $12$ and $8$, the fields of \autoref{tab:cube} change by less than $10^{-6}$, and by less than $10^{-7}$ up to $J = 10$. These node numbers are for a side of one wavelength, and they must grow with the wavenumber because the phase of the kernel varies along a panel in proportion to $k$. At two wavelengths we use $20$ and $10$ nodes; the fields then differ from those with $32$ and $14$ nodes by less than $10^{-10}$, whereas with $12$ and $8$ nodes they are wrong by up to $2\times10^{-4}$. At three wavelengths we use $28$ and $12$ nodes.

\section{Boundary element settings}
\setcounter{table}{0}
\label{app:bempp}

We used Bempp-cl version 0.4.2 \cite{betcke2021}. In this library, the electric field integral equation is discretised with lowest-order divergence-conforming Rao--Wilton--Glisson (RWG) trial functions and scaled N\'ed\'elec test functions ($\bn\times\text{RWG}$), and the resulting double-precision dense systems are solved by LU factorisation with LAPACK. The sphere is meshed by the built-in generator of the library at a prescribed element size. In contrast, the torus and the almond use structured grids in the two parameters of the surface, with $n$ cells around the tube of the torus and $2n$ around its axis, and each cell is cut into two triangles. For these three shapes we used the library's default quadrature orders, $4$ for regular and $4$ for singular integrals. Raising the orders to $4$ and $6$ or $6$ and $8$ changes the error by less than $1\%$, because the mesh sets the error.

\emph{Backend.} Bempp-cl has two assembly backends. We ran the OpenCL backend on the processor through the PoCL driver, which assembled the matrices of this paper about twenty times faster than the Numba backend, with identical results.

\emph{Quadrature and meshes for the cube.} Here we used the quadrature orders $4$ and $6$ for regular and singular integrals, $6$ and $8$, and $8$ and $10$. The graded meshes use a tensor grid on each face with the nodes $\pm\tfrac12\bigl(1-(1-|\xi|)^{g}\bigr)$, $\xi$ equally spaced in $[-1,1]$, and each rectangle is cut into two triangles. We scanned the grading exponent $g$ between $1.25$ and $2.25$ together with the three quadrature settings and seven mesh sizes with up to $28800$ unknowns, for each of the incident waves and wavenumbers that we report. On the finest mesh that fits into memory, with $41472$ unknowns, we ran the gradings from $1.5$ to $1.9$. \autoref{tab:bemenv} lists, for the standard wave at $k = 2\pi$, the runs that reduce the error by at least $3\%$ compared with every faster run. Which grading is best depends on the accuracy. In \autoref{tab:bemenv} the gradings $1.6$ to $1.9$ appear, and across all four cube cases and \autoref{fig:cubeconv} every scanned grading from $1.25$ to $2$ ($1.5$ to $2$ below $3\times10^{-3}$) appears among the best runs. Quadrature orders $4$ and $6$ are the fastest at errors above about $3.3\times10^{-4}$, whereas the raised orders are the fastest at all smaller errors in \autoref{tab:bemenv}. The Bempp curves in Figures~\ref{fig:cubecases} and~\ref{fig:costall} trace the lower envelopes formed by these runs.

\begin{table}[htbp]
  \centering
  \caption{Unit cube, $k = 2\pi$, standard wave: the fastest Bempp runs with the symmetry of the cube for each accuracy, over all gradings $g$, quadrature orders and mesh sizes that we tried. Difference from the MFS reference.}
  \label{tab:bemenv}
  \small
  \begin{tabular}{rlrrl}
    \toprule
    $g$ & Quadrature & Unknowns & Time [s] & Error \\
    \midrule
\tablebody{cube_bem_env}
    \bottomrule
  \end{tabular}
\end{table}

\emph{Symmetry.} Bempp does not use the symmetry of a body by itself. To compare the two codes on an equal footing, we therefore solved the boundary element systems with the same symmetry as the MFS: the three mirror planes of the cube, the two of the almond, and the rotational symmetry of the torus. For the mirror planes, the mesh consists of one octant or quarter of the surface and its mirror images, so that a reflection acts on the RWG coefficients as a signed permutation. A current of a given parity class is then a combination of symmetrised basis functions, one for each reflection orbit of mirror-image edges. By restricting the test space in Bempp to one octant or quarter, we assemble only the corresponding rows and sum the columns belonging to each reflection orbit with the parity signs of the class. Folding the matrix in this way lets us solve each class by an LU factorisation of an eighth or a quarter of the full size. For the cube we also use the rotation about the body diagonal as in \autoref{sec:cubepoints}, so that a general wave needs four factorisations. For the torus the rotation by one cell about the axis maps the structured mesh to itself. Since the rotations also act as signed permutations, the matrix is block-circulant, and the system splits into one small system for each Fourier order about the axis. Only the rows of one test function per orbit are assembled, and the orders that the incident wave does not excite are skipped.

The solution with symmetry differs from that of the full system by $3\times10^{-5}$ with the orders $4$ and $6$ and by $7\times10^{-6}$ with the orders $8$ and $10$, because the quadrature rules of the library are not invariant under the reflections. With the symmetry, the boundary element solution on the cube is obtained five to ten times faster, and meshes with up to $41472$ unknowns fit into memory. On the almond the gain is three to seven times, and on the torus ten to twenty times, and meshes with up to $301056$ unknowns were used. For the torus the solution with symmetry equals that of the full system to $10^{-14}$. Note that without the symmetry, the available memory on our machine limits a single dense matrix to about $4$\,GB, which is $16384$ unknowns.
